\section{Proofs omitted from the main text}\label{S:proofs}

\subsection{Conventions}
Cells $\theta\in\{\pm1\}^m$, actions $u\in\{\pm1\}^m\cup\{0\}$ (the hold is the zero vector), drift $D(u,\theta)=2\langle u,\theta\rangle-5$ in tenths, slack $s_\theta(t)=L+\sum_{\tau<t}D(u_\tau,\theta)$. A set $S$ is kept by a word from $L$ iff $s_\theta(t)\ge0$ for all $t$ and all $\theta\in S$.

\subsection{The slack game (main Proposition~\mref{prop:game})}
\begin{proof}
(a) If $x$ is winning (for $k$ steps) and $x'\ge x$ componentwise, the same word keeps $x'$, since slacks only shift upward. The recursion for $\calW_{k+1}$ is the definition; for the minimal elements note that $x\in\calW_{k+1}$ iff $x+D(u)\ge w$ for some $u$ and some $w\in\min\calW_k$, i.e. $x\ge(w-D(u))^+$, because $x\ge0$ anyway.

(b) $\calW\subseteq\bigcap_k\calW_k$ is trivial. Conversely, if $x\in\calW_k$ for every $k$, the tree of words that keep $x$ alive for their length is infinite and finitely branching ($|U_m|=2^m+1$), so it has an infinite branch by K\"onig's lemma, which is a word keeping $x$ for ever.

(c) By definition of $\Lstar$ and of blind survivability at level $L$ (all cells start at the common slack $L$).

(d) Let $x\in\calW^{c}$: a word keeps the clipped trajectory nonnegative for ever. The clipped trajectory is dominated by the unclipped one (clipping only lowers slacks, and the drifts do not depend on the state), so the same word keeps the unclipped trajectory nonnegative: $x\in\calW$. Now let $c$ dominate every coordinate of every minimal element of $\calW$ and let $x\in\calW\cap[0,c]^S$. Choose $u$ with $x+D(u)\in\calW$ (exists because $\calW$ is the greatest fixed point of the one-step operator, by (b) and the recursion in (a)); then $x+D(u)\ge w$ for some minimal $w\in\calW$, and $\min(x+D(u),c)\ge\min(w,c)=w$, so the clipped successor is again in $\calW\cap[0,c]^S$. Repeating, the player stays in $\calW\cap[0,c]^S$ for ever in the clipped game, i.e. $x\in\calW^{c}$. The computation: $\calW^{c}$ is the greatest fixed point of $X\mapsto\{x\in[0,c]^S:\exists u,\ \min(x+D(u),c)\in X\}$ on up-sets of the finite box, and the predecessor of an up-set with minimal elements $M$ has minimal elements $\min\{(w-D(u))^+:w\in M,u\}\cap[0,c]^S$: since $w_i\le c$ for every minimal $w$ of the box, $\min(x_i+D_i(u),c)\ge w_i$ iff $x_i+D_i(u)\ge w_i$, so the predecessor condition is $x\ge(w-D(u))^+$ exactly as in the unclipped game, intersected with the box. The iteration from $X_0=[0,c]^S$ is decreasing in a finite lattice and stabilises.

(e) If a word keeps $x_0$ for $t$ steps, its state after $\tau\le t$ steps is dominated by an element of $R_\tau$ (induction: the successor of a dominated state is dominated by the successor of the dominating state, which is dominated by a maximal element of $R_{\tau+1}$). So $R_t=\emptyset$ means no word survives $t$ steps, and the last nonempty $R_{t-1}$ witnesses $t-1$ steps.
\end{proof}

\subsection{The certified thresholds}
For a set $S$ the module computes $\calW^{c}(S)$ for $c=6,12,24,48,96$ in turn, records the least $L$ with $(L,\dots,L)\in\calW^{c}(S)$, and then runs the forward reachability of (e) from $(L-1,\dots,L-1)$ with a budget of $400$ steps. If the frontier empties, $\Lstar(S)=L$ is certified; otherwise the next cap is tried. For every set of this paper the first cap at which the game won was already certified except for the subsets of the boundary ten-sets, where cap $6$ returned $5$ and cap $12$ returned the certified $3$. The number of sets with an uncertified threshold is $0$ for $m=4,5,6$ (Table~S4.5).

\subsection{The sharp form of the finite-horizon corollary}
Main Corollary~\mref{cor:horizon} bounds the horizon by $\lfloor L/(5-2V)\rfloor$. The bound uses only the pooled inequality, so it is attained exactly when some word keeps the pooled slack $\sum_\theta q_\theta s_\theta(t)$ falling at the slowest rate $5-2V$ per step while no single cell of the support dies earlier. For a far pair at even $m$ the constant action at distance $\tfr{m-2}2$ from both cells gives both the drift $2m-5-4\cdot\tfr{m-2}{2}=-1$, so each cell loses exactly one tenth per step and the horizon is $L$. For $m=5$ and $h=3$ every action is at distances $(d_a,d_b)$ with $d_a+d_b\ge3$, so the drifts are $(1,-3)$, $(-3,1)$, $(-3,-3)$ or worse, and no action costs both cells exactly one tenth; the forward search gives $2\lfloor(L-1)/2\rfloor$ for $L\le20$ (Table~S2.2). For the whole cube ($V=0$) the hold attains the bound exactly.

\subsection{The feedback recursion (main Proposition~\mref{prop:feedback})}
Let $F(T,L)$ be as defined in the main text. Soundness: if $L\ge\Lstar(T)$ a blind word keeps every cell of $T$; otherwise the splitting action $u$ is played, the observed stock $L+D(u,\theta)$ identifies the class $T_k$ of the true cell, and by induction (classes are strictly smaller) a strategy keeps every cell of every class from its slack. Completeness: consider any strategy that keeps every cell of $T$ with the stock observed and let $u$ be its first action. If $u$ does not split $T$, the observation is uninformative and the strategy continues as a feedback strategy on $T$ from slack $L+D(u,\cdot)$, which is constant on $T$; if that drift is nonnegative then $T\subseteq\Ball(u,r^\ast)$ is blind-survivable at the floor, and if it is negative the continuation shows $F(T,L')$ for some $L'<L$, hence $F(T,L)$ by monotonicity. If $u$ splits $T$, the continuation is a feedback strategy on each class from its slack, so $F(T_k,L+k)$ holds for every class by induction, and $F(T,L)$ follows. The recursion is evaluated with memoisation on $(T,L)$ and the certified blind thresholds as the base case; the thresholds of the information classes met in the recursion at $m=4,5$ are all in $\{0,1,3,\infty\}$.
